% dcrgraph-cookbook.tex -- recipes for the dcrgraph package: short,
% task-oriented, each with a real graph. Compile with LuaLaTeX (the import
% recipes need it), twice, in this folder (manuals/cookbook/).
%
% The examples are in examples/: each is shown as code and executed, so the
% code shown is exactly what made the picture. In the step-by-step recipes
% every file holds the complete code up to that step; lines starting with
% \NEW are the new or changed ones: \NEW does nothing in TeX and is hidden
% in the listing, which shows those lines in green.
\documentclass[a4paper,11pt]{article}
\usepackage[margin=2.2cm]{geometry}
\usepackage[T1]{fontenc}
\usepackage{lmodern}
\usepackage[scaled=0.92]{helvet}
% dcrgraph.sty is two folders up
\makeatletter\def\input@path{{../../}{./}}\makeatother
\usepackage{dcrgraph}
\usepackage[most]{tcolorbox}
\usepackage{listings}
\usepackage[hidelinks]{hyperref}
\usepackage{parskip}
\sloppy

\colorlet{newcode}{green!45!black}
\let\NEW\relax
% long lines are wrapped, the rest of a line indented (one command per
% line in the files: \NEW must start a line)
\lstset{basicstyle=\ttfamily\footnotesize, columns=fullflexible, keepspaces,
  upquote=true, breaklines, breakindent=4em, postbreak={},
  moredelim=[il][\color{newcode}]{\\NEW}}
% \example{name}: examples/name.tex as code, and below it what it produces,
% scaled down to the width of the text if it is wider
\newsavebox\cbpicture
\newcommand\example[1]{%
  \begin{tcolorbox}[enhanced, breakable, colback=white, colframe=black!25,
      boxrule=0.4pt, arc=1mm, left=2mm, right=2mm, halign lower=center]
    \lstinputlisting{examples/#1.tex}
    \tcblower
    \sbox\cbpicture{\input{examples/#1}}%
    \ifdim\wd\cbpicture>\linewidth
      \resizebox{\linewidth}{!}{\usebox\cbpicture}%
    \else
      \usebox\cbpicture
    \fi
  \end{tcolorbox}}
\newcommand\cs[1]{\texttt{\textbackslash#1}}
\newcommand\key[1]{\texttt{#1}}
\newcommand\new{\textcolor{newcode}{green}}

\title{The \textsf{dcrgraph} cookbook\\[4pt]
  \large Recipes for DCR graphs in \LaTeX}
\author{tilzuck}
\date{Draft}

\begin{document}
\maketitle
\tableofcontents

% =====================================================================
\section*{About the recipes}
Each recipe draws a graph from the examples of the DCR Solutions
documentation, first by hand, step by step, then by including the
graph's XML file. In each step the code is complete; the lines that are
new or changed are \new. All you need is \cs{usepackage\{dcrgraph\}}; the
imports need LuaLaTeX.

% =====================================================================
\section{Events, nestings and relations}
The onboarding process from the DCR documentation,
\url{https://documentation.dcr.design/example/example-of-onboarding-process/}.

\subsection{Step 1: one event}
Everything of a graph goes into a \texttt{dcrgraph} environment. An event is
\cs{activity\{}\textit{label}\texttt{\}\{}\textit{id}\texttt{\}}; the id is the
name you refer to it by. A long label is wrapped, and the box grows if
needed.
\example{onboarding-1}

\subsection{Step 2: all events, placed next to each other}
Options go into square brackets: \key{role=} for the role in the header,
and the position relative to another event: \key{right=0.8cm of computer}
(also \key{left=}, \key{above=}, \key{below=}).
\example{onboarding-2}

\subsection{Step 3: nestings}
\cs{nesting\{}\textit{label}\texttt{\}\{}\textit{id}\texttt{\}\{}\textit{ids of its
events}\texttt{\}} is drawn around the given events, so it comes after them,
and an inner nesting before the outer one. Leave some room between events
of different nestings (here 1.8\,cm) for the margins and labels.
\example{onboarding-3}

\subsection{Step 4: conditions}
Relations are written from the source to the target,
\cs{condition\{}\textit{from}\texttt{\}\{}\textit{to}\texttt{\}}, and may start or
end at a nesting. Several relations between the same two events are
spread out side by side; this needs two \LaTeX{} runs (latexmk and Overleaf
do that by themselves).
\example{onboarding-4}

\subsection{Step 5: responses, includes and excludes}
\cs{response}, \cs{include} and \cs{exclude} work the same way. A relation
from an event to itself is drawn as a marker on its top edge.
\example{onboarding-5}

\subsection{Step 6: milestone and marking}
The marking is set with the options \key{pending}, \key{executed} and
\key{excluded} of an event.
\example{onboarding-6}

\subsection{The same graph from the editor}
The graph's XML file, exported from the DCR portal, gives the same graph in
one line, with the positions from the editor.
\example{onboarding-import}

% =====================================================================
\section{Subprocesses, subgraphs and spawns}
The subgraph example from the DCR documentation,
\url{https://documentation.dcr.design/example/subgraph/}.

\subsection{Step 1: the events}
Event types are options like the marking: here \key{computation}; also
\key{data}, \key{dmn}, \key{global}, \key{form}.
\example{subgraph-1}

\subsection{Step 2: subprocess and subgraph}
\cs{subprocess} and \cs{subgraph} are written like \cs{nesting}: label, id,
the ids of what is inside. The label may be empty.
\example{subgraph-2}

\subsection{Step 3: the relations}
The spawn goes to the subgraph, written like any other relation.
\example{subgraph-3}

\subsection{The same graph from the editor}
\example{subgraph-import}

\end{document}
